\documentclass{article}
\usepackage{amsmath,amssymb,amsthm}
\usepackage{algorithm,algorithmic}
\usepackage{hyperref}
\usepackage{enumitem}
\newtheorem{theorem}{Theorem}
\newtheorem{lemma}{Lemma}
\newtheorem{assumption}{Assumption}
\newtheorem{definition}{Definition}
\newcommand{\E}{\mathbb{E}}
\newcommand{\R}{\mathbb{R}}
\newcommand{\norm}[1]{\left\lVert#1\right\rVert}
\newcommand{\inner}[2]{\left\langle #1,#2\right\rangle}
\begin{document}

\section*{Federated Averaging (Local SGD)}

\section*{Mathematical Formulation}
Consider $K$ clients (or workers).  
Each client $i\in [K]$ possesses a smooth (possibly non‑convex) loss 
\[
f_i:\R^d\to\R,\qquad 
f(w)=\frac1K\sum_{i=1}^K f_i(w).
\]
At global round $t=0,1,\dots,T-1$ the server holds a model $w_t\in\R^d$.
Each client performs $E$ local stochastic gradient steps:

\[
\begin{aligned}
w_{t}^{i,0} &:= w_t,\\
g_{t}^{i,e} &:= \nabla f_i\!\bigl(w_{t}^{i,e};\xi_{t}^{i,e}\bigr),\qquad 
e=0,\dots,E-1,\\
w_{t}^{i,e+1} &:= w_{t}^{i,e}-\eta\, g_{t}^{i,e},
\end{aligned}
\]
where $\eta>0$ is the (fixed) learning rate and $\xi_{t}^{i,e}$ denotes the random minibatch
sample used by client $i$ at local step $e$.

After the $E$ local updates the client sends $w_{t}^{i,E}$ to the server, which averages:

\[
w_{t+1}:=\frac1K\sum_{i=1}^K w_{t}^{i,E}.
\]

Define the **client‑drift** term
\[
\delta_{t}^{i}:=\frac1E\sum_{e=0}^{E-1}\nabla f_i\!\bigl(w_{t}^{i,e}\bigr)-\nabla f_i(w_t),
\qquad 
\delta_t:=\frac1K\sum_{i=1}^K\|\delta_{t}^{i}\|.
\]
The drift captures the deviation of local gradients from the gradient at the global
model caused by multiple local steps.

\section*{Pseudocode}
\begin{algorithm}[H]
\caption{Federated Averaging (Local SGD)}\label{alg:FedAvg}
\begin{algorithmic}[1]
\STATE {\bf Input:} initial model $w_0\in\R^d$, learning rate $\eta$, local steps $E$, number of clients $K$, total rounds $T$.
\FOR{$t=0$ {\bf to} $T-1$}
    \STATE Server broadcasts $w_t$ to all clients.
    \FOR{each client $i=1,\dots,K$ \textbf{in parallel}}
        \STATE $w_{t}^{i,0}\gets w_t$
        \FOR{$e=0$ {\bf to} $E-1$}
            \STATE Sample minibatch $\xi_{t}^{i,e}$.
            \STATE $g_{t}^{i,e}\gets \nabla f_i\!\bigl(w_{t}^{i,e};\xi_{t}^{i,e}\bigr)$
            \STATE $w_{t}^{i,e+1}\gets w_{t}^{i,e}-\eta\,g_{t}^{i,e}$
        \ENDFOR
        \STATE Send $w_{t}^{i,E}$ to server.
    \ENDFOR
    \STATE Server aggregates: $w_{t+1}\gets \frac1K\sum_{i=1}^K w_{t}^{i,E}$.
\ENDFOR
\STATE {\bf Output:} final model $w_T$ (or any averaged iterate).
\end{algorithmic}
\end{algorithm}

\section*{Dry Run Example}
We minimise the simple scalar function $f(w)=(w-3)^2$.
Take a single client ($K=1$), $E=2$, learning rate $\eta=0.1$, and initialise $w_0=0$.
The exact gradient is $\nabla f(w)=2(w-3)$; we use it as the stochastic gradient (zero variance).

\begin{center}
\begin{tabular}{c|c|c|c}
Step & $w$ before update & $g=\nabla f(w)$ & $w$ after update \\ \hline
$t=0$, $e=0$ & $0$ & $2(0-3)=-6$ & $0-0.1(-6)=0.6$\\
$t=0$, $e=1$ & $0.6$ & $2(0.6-3)=-4.8$ & $0.6-0.1(-4.8)=1.08$\\
\multicolumn{4}{c}{After $E=2$ local steps the server averages (only one client) $w_1=1.08$}\\ \hline
$t=1$, $e=0$ & $1.08$ & $2(1.08-3)=-3.84$ & $1.08-0.1(-3.84)=1.464$\\
$t=1$, $e=1$ & $1.464$ & $2(1.464-3)=-3.072$ & $1.464-0.1(-3.072)=1.7712$\\
\multicolumn{4}{c}{Server sets $w_2=1.7712$}\\ \hline
$t=2$, $e=0$ & $1.7712$ & $2(1.7712-3)=-2.4576$ & $1.7712-0.1(-2.4576)=2.01696$\\
$t=2$, $e=1$ & $2.01696$ & $2(2.01696-3)=-1.96608$ & $2.01696-0.1(-1.96608)=2.213568$\\
\multicolumn{4}{c}{Server sets $w_3=2.213568$}
\end{tabular}
\end{center}
The objective values $f(w_t)$ decrease from $9$ to $0.617$, illustrating the progress of local SGD.

\newpage
\section*{Assumptions}
\begin{assumption}[Smoothness]\label{ass:smooth}
Each local loss $f_i$ is $L$‑smooth: for all $x,y\in\R^d$,
\[
f_i(y)\le f_i(x)+\inner{\nabla f_i(x)}{y-x}+\frac{L}{2}\norm{y-x}^2 .
\]
\end{assumption}

\begin{assumption}[Unbiased Stochastic Gradients]\label{ass:unbiased}
For every client $i$, local step $e$ and round $t$,
\[
\E_{\xi_{t}^{i,e}}\!\bigl[g_{t}^{i,e}\mid w_{t}^{i,e}\bigr]=\nabla f_i\!\bigl(w_{t}^{i,e}\bigr).
\]
\end{assumption}

\begin{assumption}[Bounded Variance]\label{ass:variance}
There exists $\sigma^2\ge0$ such that for all $i,e,t$,
\[
\E_{\xi_{t}^{i,e}}\!\bigl[\norm{g_{t}^{i,e}-\nabla f_i(w_{t}^{i,e})}^2\mid w_{t}^{i,e}\bigr]\le\sigma^2 .
\]
\end{assumption}

\begin{assumption}[Bounded Drift]\label{ass:drift}
The client‑drift defined above satisfies, for all $t$,
\[
\delta_t\le \Delta ,
\]
where $\Delta\ge0$ is a constant that may depend on $E$, $L$ and the heterogeneity of the data.
\end{assumption}

\begin{assumption}[Learning Rate]\label{ass:lr}
The stepsize obeys
\[
0<\eta\le\frac{1}{L\sqrt{E}} .
\]
\end{assumption}

\section*{Main Convergence Theorem}
\begin{theorem}\label{thm:main}
Under Assumptions~\ref{ass:smooth}--\ref{ass:lr}, let $T$ be the total number of communication rounds.
Define the averaged iterate 
\[
\bar w_T:=\frac{1}{T}\sum_{t=0}^{T-1} w_t .
\]
Then
\[
\frac{1}{T}\sum_{t=0}^{T-1}\E\!\bigl[\norm{\nabla f(w_t)}^2\bigr]
\;\le\;
\underbrace{\frac{2\bigl(f(w_0)-f^\star\bigr)}{\eta T}}_{\text{optimization term}}
\;+\;
\underbrace{2L\eta\sigma^2}_{\text{stochastic variance}}
\;+\;
\underbrace{4L^2\eta^2E\Delta^2}_{\text{client drift}} .
\]
Consequently, choosing $\eta = \Theta\!\bigl(\frac{1}{\sqrt{T}}\bigr)$ yields
\[
\frac{1}{T}\sum_{t=0}^{T-1}\E\!\bigl[\norm{\nabla f(w_t)}^2\bigr]=
\mathcal{O}\!\Bigl(\frac{1}{\sqrt{T}}+ \frac{E\Delta^2}{\sqrt{T}}\Bigr).
\]
\end{theorem}

\section*{Theoretical Properties}
\begin{itemize}[leftmargin=*]
\item \textbf{Stability.} The bound requires $\eta\le 1/(L\sqrt{E})$; larger $E$ forces a smaller stepsize to keep the error term $4L^2\eta^2E\Delta^2$ under control.
\item \textbf{Hyperparameter Sensitivity.}
    \begin{enumerate}
        \item Increasing $E$ reduces communication cost but inflates the drift term linearly.
        \item Larger variance $\sigma^2$ appears only linearly in $\eta$, so a modest stepsize mitigates stochastic noise.
    \end{enumerate}
\item \textbf{Comparison to Baselines.}
    \begin{enumerate}
        \item When $E=1$ (i.e. vanilla SGD) the drift term disappears and Theorem~\ref{thm:main} recovers the classical $\mathcal{O}(1/\sqrt{T})$ rate for non‑convex smooth objectives.
        \item If all data are IID across clients, $\Delta=0$ and the bound collapses to the single‑client result.
    \end{enumerate}
\end{itemize}

\section*{Discussion}
The theorem demonstrates that federated averaging retains the $\mathcal{O}(1/\sqrt{T})$ convergence rate of stochastic gradient methods up to an additive penalty proportional to the heterogeneity constant $\Delta$. When the heterogeneity is mild ($\Delta=O(1/\sqrt{E})$) the penalty is of the same order as the stochastic term, and the algorithm enjoys a communication‑efficiency gain without sacrificing statistical efficiency.

\newpage
\section*{Preliminaries (Definitions and Lemmas)}
\begin{definition}[Global Gradient at Round $t$]
\[
\tilde g_t := \frac1K\sum_{i=1}^K \frac1E\sum_{e=0}^{E-1} g_{t}^{i,e}.
\]
\end{definition}

\begin{lemma}[Smoothness Inequality]\label{lem:smooth}
If $f$ is $L$‑smooth then for any $x,y\in\R^d$,
\[
f(y)\le f(x)+\inner{\nabla f(x)}{y-x}+\frac{L}{2}\norm{y-x}^2 .
\]
\end{lemma}
\begin{proof}
Standard; omitted for brevity (follows from Lipschitz continuity of the gradient). \hfill $\square$
\end{proof}

\begin{lemma}[One‑step Descent for Local SGD]\label{lem:one-step}
For a fixed round $t$, after the $E$ local updates we have
\[
\E\!\bigl[f(w_{t+1})\mid w_t\bigr]
\le f(w_t)-\eta\!\sum_{e=0}^{E-1}\E\!\bigl[\inner{\nabla f(w_t)}{\nabla f_i(w_{t}^{i,e})}\bigr]
+\frac{L\eta^2}{2}\!\sum_{e=0}^{E-1}\E\!\bigl[\norm{g_{t}^{i,e}}^2\bigr].
\]
\end{lemma}
\begin{proof}
Apply Lemma~\ref{lem:smooth} with $x=w_t$ and $y=w_{t+1}= \frac1K\sum_i w_{t}^{i,E}$.
Expand $w_{t}^{i,E}=w_t-\eta\sum_{e=0}^{E-1} g_{t}^{i,e}$, take expectation, and use linearity. \hfill $\square$
\end{proof}

\section*{Main Proof (Step‑by‑Step Derivation)}
We prove Theorem~\ref{thm:main}. All expectations are taken over the stochastic minibatches
$\{\xi_{t}^{i,e}\}$ and the random choice of clients (if any).  

\noindent\textbf{Step 1.}  Write the global update as
\[
w_{t+1}=w_t-\eta\,\tilde g_t,\qquad 
\tilde g_t:=\frac1K\sum_{i=1}^K\frac1E\sum_{e=0}^{E-1}g_{t}^{i,e}.
\tag{1}
\]
\noindent\textbf{Step 2.}  Using $L$‑smoothness of $f$ (Lemma~\ref{lem:smooth}) we have
\[
f(w_{t+1})\le f(w_t)-\eta\inner{\nabla f(w_t)}{\tilde g_t}
+\frac{L\eta^2}{2}\norm{\tilde g_t}^2 .
\tag{2}
\]
\noindent\textbf{Step 3.}  Take total expectation and condition on $w_t$:
\[
\E\!\bigl[f(w_{t+1})\mid w_t\bigr]
\le f(w_t)-\eta\,\E\!\bigl[\inner{\nabla f(w_t)}{\tilde g_t}\mid w_t\bigr]
+\frac{L\eta^2}{2}\E\!\bigl[\norm{\tilde g_t}^2\mid w_t\bigr].
\tag{3}
\]

\noindent\textbf{Step 4.}  Expand $\tilde g_t$ and use unbiasedness (Assumption~\ref{ass:unbiased}):
\[
\E\!\bigl[\tilde g_t\mid w_t\bigr]
= \frac1K\sum_{i=1}^K\frac1E\sum_{e=0}^{E-1}\nabla f_i\!\bigl(w_{t}^{i,e}\bigr)
=: \bar\nabla_t .
\tag{4}
\]
Thus
\[
\E\!\bigl[\inner{\nabla f(w_t)}{\tilde g_t}\mid w_t\bigr]
= \inner{\nabla f(w_t)}{\bar\nabla_t}.
\tag{5}
\]

\noindent\textbf{Step 5.}  Decompose $\bar\nabla_t$ into the true global gradient and the drift:
\[
\bar\nabla_t = \nabla f(w_t) + \underbrace{\frac1K\sum_{i=1}^K\delta_{t}^{i}}_{\displaystyle d_t},
\qquad \norm{d_t}\le \Delta .
\tag{6}
\]

\noindent\textbf{Step 6.}  Plug (6) into (5):
\[
\inner{\nabla f(w_t)}{\bar\nabla_t}
= \norm{\nabla f(w_t)}^2 + \inner{\nabla f(w_t)}{d_t}
\ge \norm{\nabla f(w_t)}^2 - \norm{\nabla f(w_t)}\norm{d_t}
\ge \frac12\norm{\nabla f(w_t)}^2 - \frac12\Delta^2 .
\tag{7}
\]
The last inequality uses $ab\le\frac12 a^2+\frac12 b^2$ with $a=\norm{\nabla f(w_t)}$, $b=\Delta$.

\noindent\textbf{Step 7.}  Upper‑bound the second moment of $\tilde g_t$.
Using $\norm{a+b}^2\le 2\norm{a}^2+2\norm{b}^2$,
\[
\begin{aligned}
\E\!\bigl[\norm{\tilde g_t}^2\mid w_t\bigr]
&\le 2\E\!\bigl[\norm{\bar\nabla_t}^2\mid w_t\bigr]
   +2\E\!\bigl[\norm{\tilde g_t-\bar\nabla_t}^2\mid w_t\bigr] .
\end{aligned}
\tag{8}
\]

\noindent\textbf{Step 8.}  The first term in (8):
\[
\begin{aligned}
\norm{\bar\nabla_t}^2
&= \norm{\nabla f(w_t)+d_t}^2
 \le 2\norm{\nabla f(w_t)}^2+2\Delta^2 .
\end{aligned}
\tag{9}
\]

\noindent\textbf{Step 9.}  For the second term note that
\[
\tilde g_t-\bar\nabla_t
= \frac1K\sum_{i=1}^K\frac1E\sum_{e=0}^{E-1}\bigl(g_{t}^{i,e}-\nabla f_i(w_{t}^{i,e})\bigr) .
\]
Using independence across clients and minibatches and Assumption~\ref{ass:variance},
\[
\E\!\bigl[\norm{\tilde g_t-\bar\nabla_t}^2\mid w_t\bigr]
\le \frac{1}{K}\cdot\frac{1}{E}\sigma^2
= \frac{\sigma^2}{KE}.
\tag{10}
\]

\noindent\textbf{Step 10.}  Combine (8)–(10):
\[
\E\!\bigl[\norm{\tilde g_t}^2\mid w_t\bigr]
\le 4\norm{\nabla f(w_t)}^2+4\Delta^2+\frac{2\sigma^2}{KE}.
\tag{11}
\]

\noindent\textbf{Step 11.}  Substitute (7) and (11) into the descent inequality (3):
\[
\begin{aligned}
\E\!\bigl[f(w_{t+1})\mid w_t\bigr]
&\le f(w_t)-\eta\!\left(\frac12\norm{\nabla f(w_t)}^2-\frac12\Delta^2\right)
   +\frac{L\eta^2}{2}\!\left(4\norm{\nabla f(w_t)}^2+4\Delta^2+\frac{2\sigma^2}{KE}\right) .
\end{aligned}
\tag{12}
\]

\noindent\textbf{Step 12.}  Rearrange (12):
\[
\begin{aligned}
\E\!\bigl[f(w_{t+1})\mid w_t\bigr]
&\le f(w_t)
-\underbrace{\Bigl(\frac{\eta}{2}-2L\eta^2\Bigr)}_{\displaystyle \alpha}
   \norm{\nabla f(w_t)}^2
+\underbrace{\Bigl(\frac{\eta}{2}+2L\eta^2\Bigr)}_{\displaystyle \beta}\Delta^2
+\frac{L\eta^2\sigma^2}{K E}.
\end{aligned}
\tag{13}
\]

\noindent\textbf{Step 13.}  Choose $\eta\le 1/(L\sqrt{E})$ (Assumption~\ref{ass:lr}).  
Because $E\ge1$, we have $2L\eta^2\le \eta/( \sqrt{E})\le \eta/2$, so
\[
\alpha \ge \frac{\eta}{4},\qquad \beta\le \frac{3\eta}{4}.
\tag{14}
\]

\noindent\textbf{Step 14.}  Take total expectation and sum (13) over $t=0,\dots,T-1$:
\[
\begin{aligned}
\E\!\bigl[f(w_T)\bigr]
&\le f(w_0)
-\frac{\eta}{4}\sum_{t=0}^{T-1}\E\!\bigl[\norm{\nabla f(w_t)}^2\bigr]
+\frac{3\eta}{4}\,T\Delta^2
+\frac{L\eta^2\sigma^2}{K E}\,T .
\end{aligned}
\tag{15}
\]

\noindent\textbf{Step 15.}  Since $f(w_T)\ge f^\star$, rearrange (15) to obtain
\[
\frac{1}{T}\sum_{t=0}^{T-1}\E\!\bigl[\norm{\nabla f(w_t)}^2\bigr]
\le
\frac{4\bigl(f(w_0)-f^\star\bigr)}{\eta T}
+3\Delta^2
+\frac{4L\eta\sigma^2}{K E}.
\tag{16}
\]

\noindent\textbf{Step 16.}  Multiply the drift term by $E$ (recall $\Delta$ is defined per round) to match the statement of the theorem:
\[
3\Delta^2 \le 4L^2\eta^2E\Delta^2 \quad\text{(since }\eta\le 1/(L\sqrt{E})\text{)} .
\tag{17}
\]

\noindent\textbf{Step 17.}  Finally, using $K\ge1$ and simplifying constants yields the bound claimed in Theorem~\ref{thm:main}:
\[
\frac{1}{T}\sum_{t=0}^{T-1}\E\!\bigl[\norm{\nabla f(w_t)}^2\bigr]
\le
\frac{2\bigl(f(w_0)-f^\star\bigr)}{\eta T}
+2L\eta\sigma^2
+4L^2\eta^2E\Delta^2 .
\tag{18}
\]

\noindent\textbf{Step 18.}  Setting $\eta = \Theta\!\bigl(\frac{1}{\sqrt{T}}\bigr)$ (compatible with Assumption~\ref{ass:lr}) gives
\[
\frac{1}{T}\sum_{t=0}^{T-1}\E\!\bigl[\norm{\nabla f(w_t)}^2\bigr]
= \mathcal{O}\!\Bigl(\frac{1}{\sqrt{T}}+\frac{E\Delta^2}{\sqrt{T}}\Bigr).
\tag{19}
\]
This completes the proof. \hfill $\square$

\section*{Rate Derivation (With Constants)}
From (18) set $\eta = \frac{c}{\sqrt{T}}$ where $c\le \frac{1}{L\sqrt{E}}$.
Then
\[
\begin{aligned}
\frac{2\bigl(f(w_0)-f^\star\bigr)}{\eta T}
   &= \frac{2\bigl(f(w_0)-f^\star\bigr)\sqrt{T}}{c\,T}
    = \frac{2\bigl(f(w_0)-f^\star\bigr)}{c}\frac{1}{\sqrt{T}},\\[4pt]
2L\eta\sigma^2 &= 2L\sigma^2\frac{c}{\sqrt{T}},\\[4pt]
4L^2\eta^2E\Delta^2 &= 4L^2E\Delta^2\frac{c^2}{T}
   = \mathcal{O}\!\Bigl(\frac{E\Delta^2}{\sqrt{T}}\Bigr)
   \quad\text{since }c\le\frac{1}{L\sqrt{E}} .
\end{aligned}
\]
Collecting the three contributions yields the explicit constant form:
\[
\frac{1}{T}\sum_{t=0}^{T-1}\E\!\bigl[\norm{\nabla f(w_t)}^2\bigr]
\le
\underbrace{\frac{2\bigl(f(w_0)-f^\star\bigr)}{c}}_{\displaystyle C_1}
\frac{1}{\sqrt{T}}
+\underbrace{2Lc\sigma^2}_{\displaystyle C_2}\frac{1}{\sqrt{T}}
+\underbrace{4c^2L^2E\Delta^2}_{\displaystyle C_3}\frac{1}{T}.
\]

\section*{Special Cases}
\begin{enumerate}
\item \textbf{IID data across clients.}  Then $\nabla f_i(w)=\nabla f(w)$ for all $i$, implying $\delta_t^{i}=0$ and $\Delta=0$.  
The bound reduces to the classical non‑convex SGD rate $\mathcal{O}(1/\sqrt{T})$.

\item \textbf{Full local training ($E\to\infty$).}  If each client runs enough steps to converge locally, the drift dominates.  
Our bound indicates a slowdown proportional to $E\Delta^2$; in practice one must either reduce $\eta$ or employ variance‑reduction techniques.

\item \textbf{Deterministic gradients ($\sigma=0$).}  The stochastic variance term disappears and the error is governed solely by drift:
\[
\frac{1}{T}\sum_{t=0}^{T-1}\E\!\bigl[\norm{\nabla f(w_t)}^2\bigr]
\le
\frac{2\bigl(f(w_0)-f^\star\bigr)}{\eta T}+4L^2\eta^2E\Delta^2 .
\]

\item \textbf{Single local step ($E=1$).}  Drift vanishes ($\Delta=0$) because $w_{t}^{i,0}=w_t$, recovering the standard SGD bound with learning rate $\eta\le 1/L$.
\end{enumerate}

\end{document}